\begin{scanpage}{0041}{24}
In the case when both faces of the wedge are illuminated, the field scattered by the second face also is determined by Equations (3.27) and (3.28) in which one need only replace $\phi_0$ by $\alpha-\phi_0$, and $\phi$ by $\alpha-\phi$.

Then adding the field radiated by the nonuniform part of the current with the incident wave (1.24), we find the diffraction field in the physical optics approach. This equals
\[\frac{H_z}{H_{0z}}=\begin{cases}v_2^+(\phi,\phi_0)+e^{-ikr\cos(\phi-\phi_0)}+e^{-ikr\cos(\phi+\phi_0)}\\\hfill\text{with }0<\phi<\pi-\phi_0,\\v_2^+(\phi,\phi_0)+e^{-ikr\cos(\phi-\phi_0)}\\\hfill\text{with }\pi-\phi_0<\phi<\pi,\\v_2^-(\phi,\phi_0)+e^{-ikr\cos(\phi-\phi_0)}\\\hfill\text{with }\pi<\phi<\pi+\phi_0,\\v_2^-(\phi,\phi_0)\hfill\text{with }\pi+\phi_0<\phi<\alpha,\end{cases}\tag{3.29}\]
if one face of the wedge is illuminated, and
\[\frac{H_z}{H_{0z}}=\begin{cases}v_2^+(\phi,\phi_0)+v_2^-(\alpha-\phi,\alpha-\phi_0)+\\\qquad+e^{-ikr\cos(\phi-\phi_0)}+e^{-ikr\cos(\phi+\phi_0)}\\\hfill\text{with }0<\phi<\pi-\phi_0,\\v_2^+(\phi,\phi_0)+v_2^-(\alpha-\phi,\alpha-\phi_0)+\\\qquad+e^{-ikr\cos(\phi-\phi_0)}\quad\text{with }\pi-\phi_0<\phi<\alpha-\pi,\\v_2^+(\phi,\phi_0)+v_2^+(\alpha-\phi,\alpha-\phi_0)+\\\qquad+e^{-ikr\cos(\phi-\phi_0)}\quad\text{with }\alpha-\pi<\phi<\pi,\\v_2^-(\phi,\phi_0)+v_2^+(\alpha-\phi,\alpha-\phi_0)+\\\qquad+e^{-ikr\cos(\phi-\phi_0)}\quad\text{with }\pi<\phi<2\alpha-\pi-\phi_0,\\v_2^-(\phi,\phi_0)+v_2^+(\alpha-\phi,\alpha-\phi_0)+\\\qquad+e^{-ikr\cos(\phi-\phi_0)}+e^{-ikr\cos(2\alpha-\phi-\phi_0)}\\\hfill\text{with }2\alpha-\pi-\phi_0<\phi<\alpha,\end{cases}\tag{3.30}\]
if both of its faces are illuminated.
\end{scanpage}
